\documentclass[11pt]{article}
\usepackage[margin=1in]{geometry}
\usepackage{amsmath,amssymb,amsthm,booktabs,array}
\usepackage[T1]{fontenc}
\usepackage[colorlinks=true,linkcolor=blue,urlcolor=blue,citecolor=blue]{hyperref}
\hypersetup{pdftitle={Intrinsic interval consequences of Gern's type D thesis},
pdfauthor={Research draft prepared with an Astra research team}}
\newtheorem{theorem}{Theorem}[section]
\newtheorem{corollary}[theorem]{Corollary}
\newtheorem{proposition}[theorem]{Proposition}
\newtheorem{lemma}[theorem]{Lemma}
\theoremstyle{definition}\newtheorem{definition}[theorem]{Definition}
\theoremstyle{remark}\newtheorem{remark}[theorem]{Remark}
\newcommand{\ellW}{\ell}
\newcommand{\FC}{\mathrm{FC}}
\newcommand{\CD}{\mathcal C_D}
\title{Intrinsic interval consequences of Gern's type $D$ thesis\\
\large A research note using combinatorial invariance}
\author{Research draft prepared with an Astra research team}
\date{October 6, 2026}
\begin{document}
\maketitle
\begin{abstract}
We extract consequences of Gern's thesis that apply to arbitrary equal-parameter
Coxeter systems by using Theorem 1.1 of the September 24, 2026 combinatorial
invariance preprint. A type $D$ interval with fully commutative bottom becomes
a certificate for a $0$--$1$ leading coefficient in every realization of its
abstract interval poset. More specifically, Gern's bad-element intervals give
an infinite family of degree bounds and exact leading coefficients, including
vanishing of increasingly many top coefficients. All these statements are
conditional on the cited invariance theorem; we do not independently certify
its full proof. We distinguish these formal consequences from previously
known type $B$ bounds, parabolic-coset transfers, and subsystem pattern
embeddings. No new whole-type theorem or publication-level novelty claim is
established here.
\end{abstract}

\section{Inputs and the precise scope of transfer}
Throughout, $P_{u,v}(q)$ is the ordinary equal-parameter Kazhdan--Lusztig
polynomial of a Coxeter system $(W,S)$. We use
\[
 P_{u,u}=1,\qquad P_{u,v}=0\ (u\not\leq v),\qquad
 \deg P_{u,v}\leq\left\lfloor\frac{\ell(v)-\ell(u)-1}{2}\right\rfloor
 \quad(u<v).
\]
For $u<v$ put $d=\ell(v)-\ell(u)$ and define
\[
 \mu(u,v)=\begin{cases}
 [q^{(d-1)/2}]P_{u,v}(q),&d\text{ odd},\\
 0,&d\text{ even}.
 \end{cases}
\]
For odd $d$, $\mu$ is the coefficient at the largest \emph{allowed}
degree; it need not be the leading nonzero coefficient. For even
$d$, $\mu$ is zero by convention.

The two external inputs are:
\begin{enumerate}
\item Gern \cite[Theorem 4.5.11]{gern}: in finite type $D_n$, if $x$ is
fully commutative and $x<w$, then $\mu(x,w)\in\{0,1\}$.
\item OpenAI \cite[Theorem 1.1]{ci}: for arbitrary Coxeter systems,
an order isomorphism $[u,v]\cong[x,w]$ preserves the ordinary
equal-parameter polynomial: $P_{u,v}=P_{x,w}$.
\end{enumerate}
The second statement includes infinite and noncrystallographic systems.
It requires an isomorphism of the \emph{whole closed interval}.
An embedding, equal rank numbers, an undirected graph match, or a
reflection-subgroup inclusion is not sufficient.
The consequences proved below invoke this input, rather than silently
treating the preprint as an independently verified theorem.

\begin{definition}
Let $\CD$ be the class of finite posets isomorphic to an interval $[x,w]$
in some $D_n$, $n\geq4$, with $x$ fully commutative.
A witness consists of $n,x,w$ and an explicit order isomorphism.
\end{definition}

\begin{theorem}[Type $D$ witness certificate]\label{thm:transfer}
Assume the invariance input. In any Coxeter system, if $[u,v]\in\CD$
and $u<v$, then $\mu(u,v)\in\{0,1\}$.
The element $u$ is not required to be fully commutative.
\end{theorem}
\begin{proof}
Choose a witness $\iota:[u,v]\longrightarrow[x,w]$.
Bruhat intervals are finite graded posets, and an order isomorphism
preserves their rank. Hence
$\ell(v)-\ell(u)=\ell(w)-\ell(x)$.
The invariance input gives $P_{u,v}=P_{x,w}$. For odd rank the coefficients
defining $\mu$ are consequently equal; for even rank both are zero.
Gern's theorem applies to the witness bottom $x$.
\end{proof}

\begin{corollary}[Obstruction]\label{cor:obstruction}
Assume the invariance input. If $\mu(u,v)\geq2$, then $[u,v]$ has no
fully commutative-bottom type $D$ witness.
\end{corollary}
\begin{proof}This is the contrapositive of Theorem~\ref{thm:transfer}.\end{proof}

This certificate theorem is an immediate consequence of the two inputs,
not a classification of $\CD$ or an independent new mechanism.
In particular, it does not prove the fully commutative-bottom conjecture
throughout an exceptional or affine type: that would require witnesses
for \emph{all} the relevant intervals or a separate argument.

\begin{proposition}[Fully commutative status is not encoded across types]
The abstract principal lower interval does not determine whether its top
is fully commutative across arbitrary Coxeter systems.
\end{proposition}
\begin{proof}
Take $s,t$ with $m(s,t)=3$, and separately $s',t'$ with $m(s',t')=4$.
The intervals $[e,sts]$ and $[e,s't's']$ each have levels
\[
 \{e\},\quad\{s,t\},\quad\{st,ts\},\quad\{sts\},
\]
with both rank-one elements below both rank-two elements (and the
corresponding primed levels in the second interval).
This follows directly from the subword criterion. They are therefore
isomorphic six-element posets. In $I_2(3)$ the braid $sts=tst$ makes
$sts$ non-fully-commutative. In $I_2(4)$ a reduced word of length three
cannot contain the noncommuting braid of length four; $s't's'$ is fully
commutative. See also Stembridge's definition and braid characterization
\cite{stembridge}.
\end{proof}

\section{An infinite family carrying the thesis's exact calculations}
We use Gern's numbering: $s_1$ sends the first two signed entries
$(a,b)$ to $(-b,-a)$, $s_2$ interchanges the first two entries, and
$s_i$ interchanges positions $i-1,i$ for $i\geq3$.
Right multiplication has these effects on signed one-line notation.
For even $n\geq4$ define
\[
 x_n=s_1s_2s_4s_6\cdots s_n,
 \qquad
 w_n(i)=\begin{cases}
 (-1)^{n/2},&i=1,\\
 i,&i>1\text{ odd},\\
 -(n+2-i),&i\text{ even}.
 \end{cases}
\]
Set $I_n=[x_n,w_n]$. These are Gern's Definition 2.3.8 and
Corollary 2.2.19. The factors of $x_n$ commute, so $x_n$ is fully
commutative. The formulas also give an even number of negative entries
in $w_n$, as required for $D_n$.
For odd $n\geq5$, the corresponding thesis elements are the even
$(n-1)$ elements with an additional fixed last entry; their intervals
are identical after standard parabolic inclusion.

The following exact inputs are from the thesis:
\begin{align}
 L_n:=\ell(w_n)&=\frac{3n^2}{8}+\frac n4,
 &\ell(x_n)&=\frac n2+1,\label{eq:lengths}\\
 d_n:=\operatorname{rank} I_n&=\frac{3n^2}{8}-\frac n4-1,\label{eq:rank}\\
 a_n:=a(w_n)&=\begin{cases}
 3n/4,&n\equiv0\pmod4,\\
 (3n+2)/4,&n\equiv2\pmod4,
 \end{cases}\label{eq:a}\\
 P_{x_n,w_n}&=P_{e,w_n},\label{eq:equality}\\
 \mu(x_n,w_n)&=\begin{cases}1,&n=6,\\0,&n\ne6.\end{cases}\label{eq:mu}
\end{align}
The locators are Lemma 2.2.20, Definition 2.3.8, Proposition 4.5.1,
Lemma 2.3.9, and Lemmas 4.5.2--4.5.7, respectively.
Equation~\eqref{eq:mu} uses the explicit computer calculations at
$n=6$ and $n=8$ recorded in Lemma 4.5.5; it is not entirely a symbolic
argument in the original thesis.
The thesis's Proposition 1.3.5, citing Lusztig \cite[1.3(a)]{lusztig}, gives
\[
 \deg P_{e,w_n}\leq \frac{L_n-a_n}{2}.
\]

\begin{theorem}[Bad-interval degree and coefficient transfer]\label{thm:bad}
Assume the invariance input. Let $[u,v]$ be an interval in any Coxeter
system, and suppose $[u,v]\cong I_n$ for even $n\geq4$. Put
\[
 M_n=\left\lfloor\frac{d_n-1}{2}\right\rfloor,
 \qquad A_n=\left\lfloor\frac{L_n-a_n}{2}\right\rfloor,
 \qquad B_n=\min\{M_n,A_n\}.
\]
Then
\[
 \mu(u,v)=\begin{cases}1,&n=6,\\0,&n\ne6,\end{cases}
 \qquad \deg P_{u,v}\leq B_n.
\]
For $n=8$ the bound improves to $\deg P_{u,v}\leq9$.
\end{theorem}
\begin{proof}
Rank preservation identifies the rank with $d_n$, and invariance gives
$P_{u,v}=P_{x_n,w_n}=P_{e,w_n}$. The exact $\mu$ assertion follows from
\eqref{eq:mu}. The ordinary degree bound gives $M_n$, and Lusztig's
bound with \eqref{eq:a} gives $A_n$, taking floors because polynomial
degrees are integral. When $n=8$, $d_8=21$ and $M_8=A_8=10$.
Equation~\eqref{eq:mu} says the coefficient of $q^{10}$ vanishes.
Thus the degree is at most nine.
\end{proof}

\begin{corollary}[Increasing gaps at the allowed top degree]\label{cor:gaps}
Assume the invariance input and $[u,v]\cong I_n$.
If $n=8k$ with $k\geq2$, then
\[
 [q^{M_n-j}]P_{u,v}=0\quad(0\leq j\leq k-2).
\]
If $n=8k+6$ with $k\geq1$, then
\[
 [q^{M_n-j}]P_{u,v}=0\quad(0\leq j\leq k-1).
\]
These assertions concern the top $k-1$ and $k$ \emph{allowed}
coefficients, respectively, without requiring any target cell or
target $a$-function calculation.
\end{corollary}
\begin{proof}
For $n=8k$ one obtains
\[
 M_n=12k^2-k-1,\qquad A_n=12k^2-2k,
 \qquad M_n-A_n=k-1.
\]
For $n=8k+6$ one obtains
\[
 M_n=12k^2+17k+5,\qquad A_n=12k^2+16k+5,
 \qquad M_n-A_n=k.
\]
Every exponent listed in the corollary is strictly greater than $A_n$;
Theorem~\ref{thm:bad} therefore forces its coefficient to vanish.
\end{proof}

\begin{center}
\begin{tabular}{rrrrrr}
\toprule
$n$&$L_n$&$\ell(x_n)$&$d_n$&degree bound&$\mu$\\
\midrule
4&7&3&4&1&0\\
6&15&4&11&5&1\\
8&26&5&21&9&0\\
10&40&6&34&16&0\\
12&57&7&50&24&0\\
14&77&8&69&33&0\\
16&100&9&91&44&0\\
22&187&12&175&85&0\\
24&222&13&209&102&0\\
\bottomrule
\end{tabular}
\end{center}
This table is arithmetic derived from the thesis formulas; it does not
claim to recompute the high-rank polynomials. The special row $n=8$
includes the improvement from Lemma 4.5.5.

\begin{proposition}[Exact rank-eleven polynomial; computer assisted]\label{prop:d6}
The project's exact recurrence computation gives
\[
 F(q):=P^{D_6}_{x_6,w_6}(q)=1+6q+11q^2+6q^3+q^4+q^5.
\]
Assuming the invariance input, every realization $[u,v]\cong I_6$
has $P_{u,v}=F$ and $P_{u,v}(1)=26$.
\end{proposition}
\begin{proof}[Computational proof and transfer]
The endpoints, with lengths four and fifteen, are
\[
 x_6=(-1,-2,4,3,6,5),\qquad w_6=(-1,-6,3,-4,5,-2).
\]
The saved Python implementation uses signed-inversion lengths,
the Bruhat lifting property and the exact integer KL recurrence.
Its separate $R$-polynomial recurrence checks
\[
 q^{11}P_{x_6,w_6}(q^{-1})=
 \sum_{x_6\leq z\leq w_6}R_{x_6,z}(q)P_{z,w_6}(q).
\]
The closed interval has 1676 vertices and rank vector
\[
 (1,12,55,140,248,339,360,287,162,59,12,1).
\]
The lower ideal $[e,w_6]$ has 3184 vertices.
The computation also checks each right-descent choice at $w_6$;
its detailed recurrence certificate and verifier are saved with the
project. The verifier never calls the KL evaluator: it checks 24,245
stored dependency records, including 17,148 strict recurrence
identities and 6978 incomparable zero cases. All recurrence
dependencies have shorter upper length, so these identities with
the diagonal and zero base cases certify the answer by induction.
This is a computer-assisted finite calculation, separate
from the thesis's assertion of $\mu=1$.
The final transfer follows from the external invariance input.
\end{proof}

\begin{theorem}[Coefficients in abstract products]\label{thm:product}
Assume the invariance input. Suppose $[u,v]$ is isomorphic to
$\prod_{i=1}^k[x_i,w_i]$, where each factor is a nontrivial type
$D$ interval with fully commutative $x_i$. Write
$d_i=\ell(w_i)-\ell(x_i)$ and $d=\sum_i d_i$. Then
\[
 P_{u,v}(q)=\prod_i P_{x_i,w_i}(q),\qquad
 \deg P_{u,v}\leq\sum_i\left\lfloor\frac{d_i-1}{2}\right\rfloor.
\]
For $k\geq2$ one has $\mu(u,v)=0$.
If every $d_i$ is odd, the coefficient of $q^{(d-k)/2}$ is in
$\{0,1\}$. In particular, if $[u,v]\cong I_6^k$, then
$P_{u,v}=F^k$, its degree is $5k$, and its actual leading
coefficient is one.
\end{theorem}
\begin{proof}
The product poset is a Bruhat interval in the product Coxeter group.
The Hecke algebra is the tensor product of the factor Hecke algebras;
the tensor products of the canonical basis elements are bar invariant
with the required triangular expansion, hence form the canonical basis.
Their expansion gives the stated factorization of $P$.
Invariance transfers it to $[u,v]$.
The degree bound follows by adding factor bounds. Its right side is
at most $(d-k)/2<(d-1)/2$ for $k\geq2$, which implies the $\mu$ assertion.
If all $d_i$ are odd, reaching exponent $(d-k)/2$ requires the largest
allowed exponent in every factor. That coefficient is
$\prod_i\mu(x_i,w_i)\in\{0,1\}$ by Gern's theorem.
Finally $F$ has degree five and leading coefficient one.
\end{proof}
Tensor-product factorization is classical. The unrestricted transfer
from an abstract product decomposition is the use of full invariance;
no independent novelty claim is made for the product statement.

\section{Certified small examples in types \texorpdfstring{$B$ and $A$}{B and A}}
Use generators $b_0,b_1,b_2$ of $B_3$, with $b_0$ negating the
first signed entry, $b_1$ swapping the first two, and $b_2$ swapping
the last two. The type $D_4$ interval
\[
 [(1,3,2,4),(-1,2,-3,4)]
 =[s_3,s_3s_1s_2s_3]
\]
is isomorphic to the type $B_3$ interval
\[
 [(3,2,-1),(2,-1,-3)].
\]
Their endpoint lengths are $(1,4)$ and $(4,7)$, respectively;
both have rank vector $(1,4,4,1)$ and polynomial $1+q$.
The target bottom has reduced word $b_0b_1b_2b_1$ and is not
fully commutative, because $b_1b_2b_1$ is a noncommuting braid.

Here is the complete vertex bijection found by the computation:
\begin{center}\small
\begin{tabular}{rrr}
\toprule
relative rank&$D_4$ vertex&$B_3$ vertex\\
\midrule
0&$(1,3,2,4)$&$(3,2,-1)$\\
1&$(-2,3,-1,4)$&$(3,-2,-1)$\\
1&$(2,3,1,4)$&$(3,1,-2)$\\
1&$(-3,-1,2,4)$&$(-3,2,-1)$\\
1&$(3,1,2,4)$&$(1,2,-3)$\\
2&$(-1,3,-2,4)$&$(3,-1,-2)$\\
2&$(-3,2,-1,4)$&$(2,-3,-1)$\\
2&$(3,2,1,4)$&$(2,1,-3)$\\
2&$(-1,-3,2,4)$&$(-1,2,-3)$\\
3&$(-1,2,-3,4)$&$(2,-1,-3)$\\
\bottomrule
\end{tabular}\end{center}
The verifier checks all pairwise order relations, not just rank numbers.
It also checks every subinterval's KL polynomial, $R$-reciprocity,
and every available right descent in the KL recurrence.

The $B_3$ interval is not contained in a left coset of a type $A$
parabolic, even a conjugate one: two of its right cover reflections
are $(-3,2,-1)$ and $(-1,2,3)$, whose product has order four.
Two distinct reflections in a type $A$ Coxeter group have product
order two or three; the cover labels of a coset of such a subgroup
would lie in that subgroup. A corresponding pair of left cover labels
of order four excludes right type $A$ parabolic cosets as well.
Thus the target example actually encounters the double-bond structure,
although its KL equality is already covered by short-interval results.

Four additional saved certificates include $B_4$ intervals with non-fully-
commutative bottom and ranks three and five, as well as type $A_4$
examples of those ranks. The rank-five examples among these four have $P=1+q$ and
$\mu=0$, illustrating the distinction between actual leading
coefficient and $\mu$. These small examples are verification
demonstrations, not new results in those types.

A further rank-five certificate uses all four $D_4$ simple generators:
\[
 [(-2,-1,3,4),(-4,-3,-2,-1)]_{D_4}
 \cong [(2,4,3,-1),(2,3,-1,-4)]_{B_4}.
\]
The endpoint lengths are $(1,6)$ and $(5,10)$. Each interval has
34 vertices. Its rank vector and polynomial are
\[
 (1,4,10,12,6,1),\qquad P=1+2q+q^2,\qquad \mu=1.
\]
The $D_4$ bottom is $s_1$, and a reduced word of the top is
$s_1s_3s_2s_4s_3s_1$. The $B_4$ bottom has reduced word
$b_0b_1b_2b_3b_2$, containing the braid $b_2b_3b_2$.
The source's top has full $D_4$ support, and the target also has
saved order-four cover-reflection obstructions to type $A$ cosets.
This is a useful example with a fork in the source and double bond
in the target; its rank remains within previously known invariance.

\section{A structural obstruction and the shape of \texorpdfstring{$I_6$}{I6}}
\begin{proposition}[Eight-element dihedral obstruction]\label{prop:butterfly}
Let $K=[e,w_0]$ be the principal interval in $I_2(4)=B_2$,
with level sizes $(1,2,2,2,1)$. No interval in a finite simply
laced Weyl group is isomorphic to $K$. Consequently, any interval
containing a closed subinterval isomorphic to $K$ has no finite
type $D$ realization, regardless of its bottom's FC status.
\end{proposition}
\begin{proof}
We give the classical reflection-plane argument; see Dyer
\cite[Lemma 3.1 and Section 3.4]{dyer} for its antecedents.
Label a cover $z\lessdot z'$ by the reflection $z'z^{-1}$ and its
positive root. In a diamond with bottom $z$, middle vertices
$r_\alpha z,r_\beta z$, and top
$r_\rho r_\alpha z=r_\sigma r_\beta z$, one has
\[
 r_\rho r_\alpha=r_\sigma r_\beta.
\]
In a finite real reflection representation, a nonidentity product
of two distinct reflections has
$\operatorname{im}(r_\rho r_\alpha-1)=\operatorname{span}(\rho,\alpha)$.
The same image is $\operatorname{span}(\sigma,\beta)$.
The first two roots $\alpha,\beta$ are distinct and independent, so
all four diamond labels lie in their common plane.

In $K$, every vertex at one of the three middle levels is covered
by both vertices at the next middle level, when that level exists.
Apply the diamond argument first from the bottom to each rank-two
vertex. All labels at ranks zero through two lie in one plane.
Apply it next to diamonds from rank one to rank three, and then
from rank two to the top. At each step the two already known
outgoing reflections are distinct, so their roots span that same
plane. All cover labels consequently lie in that plane.

The eight vertices must then lie in one coset of the reflection
subgroup generated by roots in this plane: reach each vertex by a
cover path from the bottom. In a finite simply laced root system,
such a rank-two subsystem is $A_1\times A_1$ or $A_2$, whose
reflection groups have four or six elements. A coset cannot contain
eight distinct vertices, a contradiction.
The last assertion follows by restricting any proposed interval
isomorphism to the specified closed subinterval.
\end{proof}
Since $e$ is FC in $B_2$ and $P_{e,w_0}=1$, this obstruction is
stronger than a test using $\mu\geq2$. In particular, it disproves
the proposed shortcut that every FC-bottom interval in every finite
type can be modeled by a finite-$D$ FC-bottom interval. It is a
classical structural obstruction, not a new result of this project.

\begin{proposition}[Shape certificates for $I_6$; computer assisted]
The interval $I_6$ is not isomorphic to a principal lower interval
in any Coxeter system, and it is not a lattice.
\end{proposition}
\begin{proof}
The saved Bruhat-order computation gives rank eleven and twelve
atoms. A principal interval $[e,z]$ of rank eleven has at most eleven
atoms: they are the distinct simple generators occurring in a
reduced expression of $z$, which has eleven letters.
This proves the first assertion.

For the second, take the following two atoms of $I_6$:
\[
 a=(-4,-2,1,3,6,5),\qquad b=(-1,-3,4,2,6,5).
\]
Their two distinct minimal common upper bounds are
\[
 c=(-4,-3,1,2,6,5),\qquad d=(-2,-4,3,1,6,5).
\]
They have relative ranks two and three, respectively, and are
incomparable. The saved script enumerates all 720 common upper
bounds, verifies $c,d$ are exactly the minimal ones, and checks
the comparisons using the Bruhat lifting property. Thus $a,b$
have no join, so $I_6$ is not a lattice. This order-only certificate
does not use a KL-polynomial calculation.
\end{proof}

\section{Concrete other-type families and why they are classical}
The transfer hypotheses can be met in other types without any
new invariance theorem. The following explains both a useful application
and its novelty limitation.

\begin{lemma}[Convex parabolic cosets]\label{lem:coset}
Let $J\subseteq S$, let $W_J$ be its standard parabolic subgroup, and
let $a\in W^J$ be the minimal element of $aW_J$.
For $x\leq w$ in $W_J$, multiplication by $a$ identifies
$[x,w]_{W_J}$ with the entire interval $[ax,aw]_W$.
\end{lemma}
\begin{proof}
Parabolic decomposition $z=z^Jz_J$ is length-additive, and the
projection $z\mapsto z^J$ is order preserving. If $ax\leq z\leq aw$,
then $a\leq z^J\leq a$, hence $z\in aW_J$.
On this coset, $az_1\leq az_2$ if and only if $z_1\leq z_2$ in $W_J$.
These are the standard parabolic Bruhat-order facts; see
\cite[Section 2.5]{bb}. They identify precisely the vertices and order
of the two intervals.
\end{proof}

\begin{proposition}[Classical coset polynomial equality]\label{prop:coset}
Under the hypotheses of Lemma~\ref{lem:coset}, the equality
\[
 P^W_{ax,aw}=P^{W_J}_{x,w}
\]
holds independently of the invariance input.
\end{proposition}
\begin{proof}
Induct on $\ell(w)$. The case $x=w$ is immediate. Choose a right
descent $s\in J$ of $w$. Length additivity shows that the right descents
in $J$ of $ax$ and $aw$ are the corresponding descents of $x$ and $w$.
Use the standard right-handed KL recurrence with $s$.
Its first two terms have upper element $aws=a(ws)$, so the induction
hypothesis applies (and incomparable terms are zero).
Every nonzero summand in its correction sum has
$ax\leq z\leq a(ws)$ and $zs<z$. Lemma~\ref{lem:coset} puts $z=ay$
with $x\leq y\leq ws$ in $W_J$.
The descent condition and all length differences match. The polynomial
and $\mu$ factors in that sum have smaller upper parabolic length and
match by induction. Thus this recurrence is exactly the recurrence
in $W_J$. The recurrence is Gern \cite[Proposition 1.2.6]{gern},
with left and right exchanged, or \cite[Section 6.1]{bb}.
\end{proof}

For example, take an irreducible affine Weyl group $\widetilde D_n$ and
its finite standard $D_n$ parabolic $J$. With $\rho^\vee$ half the sum
of the positive coroots, let
\[
 a_m=t_{-2m\rho^\vee},\qquad m\geq1.
\]
Since $2\rho^\vee$ lies in the coroot lattice, these are affine Weyl
group elements. Use the affine-root convention
$t_\lambda(\alpha+r\delta)=\alpha+(r-\langle\lambda,\alpha\rangle)\delta$.
For finite simple $\alpha_i$, $a_m(\alpha_i)=\alpha_i+2m\delta$ is positive;
thus $a_m\in W^J$. Conversely,
$a_m^{-1}(\alpha_i)=\alpha_i-2m\delta$ is negative, so every finite
simple generator lies in the left descent set of $a_m$.
An adjacent pair of finite simple generators consequently forces a
reduced prefix $sts$, with $m(s,t)=3$, in $a_m$ by parabolic
decomposition. Therefore $a_m$ and its length-additive extensions
$a_m x_n$ are non-fully-commutative. This descent argument is also
Gern \cite[Lemma 1.1.21 and Corollary 1.1.22]{gern}.
Affine translations and the root-sign length criterion are standard;
see \cite[Chapters 4 and 5]{humphreys}.

Lemma~\ref{lem:coset} and Proposition~\ref{prop:coset} now give
\[
 [a_mx_n,a_mw_n]\cong I_n,
 \quad P^{\widetilde D_n}_{a_mx_n,a_mw_n}=P^{D_n}_{x_n,w_n}.
\]
In particular, $\mu=1$ for $n=6$ and $\mu=0$ for the other even
values, with all the degree bounds in Theorem~\ref{thm:bad}.
The endpoint lengths tend to infinity as $m$ tends to infinity.
This gives explicit affine examples with non-fully-commutative bottom,
but Proposition~\ref{prop:coset} makes their equality classical.
They cannot be presented as evidence that full invariance was needed.

Finite $E_7$ and $E_8$ contain standard $D_6$ parabolics, and $E_6$
contains a $D_4$ parabolic. Their standard parabolic intervals, and
coset copies as above, provide analogous exceptional-type examples.
Again these are already covered by parabolic theory, and do not
settle the fully commutative-bottom bound for all of an exceptional type.

For fully explicit rank-eleven exceptional examples, use the numbering
$1$--$3$--$4$--$5$--$6$--$7$ (and $7$--$8$ in $E_8$), with
$2$ attached to $4$. Set $J=\{2,3,4,5,6,7\}\cong D_6$ and relabel
Gern's generators $(1,2,3,4,5,6)$ as $(2,3,4,5,6,7)$.
Let $x,w$ be the resulting $x_6,w_6$, and put
$a=w_0(E_r)w_0(J)$ for $r=7,8$. This is the minimal element of
$aW_J$, of length $33$ in $E_7$ and $90$ in $E_8$.
Consequently $(ax,aw)$ has lengths $(37,48)$ and $(94,105)$,
respectively, and polynomial $F$ from Proposition~\ref{prop:d6}.
The saved exact integer-root script checks that both lower endpoints
are non-fully-commutative: stripping a displayed sequence of right
descents from $a$ leaves two adjacent descents, hence a reduced
noncommuting braid. Concatenation with a reduced word for $x$ preserves
that braid. The script and application notes contain the full sequences
and root images. These particular pairs are classical coset applications.

\section{Literature and novelty audit}
The audit is a documented search through October 6, 2026, not a proof
of the absence of every unpublished or inaccessible result.
\begin{center}
\small
\begin{tabular}{p{.35\textwidth}p{.56\textwidth}}
\toprule
Candidate statement&Assessment and primary reference\\
\midrule
Finite type $A$ FC-bottom $0$--$1$ bound&Known: Green
\cite[Theorem 4.1; arXiv Theorem 3.1]{green}. His result also covers affine type $A$.\\
Finite type $D$ FC-bottom bound&The supplied thesis:
Gern \cite[Theorem 4.5.11]{gern}.\\
Finite type $B/C$ FC-bottom bound&Known: Chmutov's 2014 thesis
\cite[Theorem 2.3.8]{chmutov}. Types $B$ and $C$ have the same
ordinary equal-parameter Coxeter system.\\
Both endpoints FC in $B$ and $H$&Known: Green
\cite[Remark 7.13]{jonestraces}. This does not alone prove the
FC-bottom/arbitrary-upper statement.\\
Upper endpoint Deodhar&Known: Jones \cite[Theorem 1.2]{jones},
including exceptional finite Weyl types. This uses a different
endpoint hypothesis; the bad elements $w_n$ are not FC.\\
Subsystem interval pattern transfer&Known: Woo
\cite[Corollary 3.3]{woo}; the arXiv numbering is Corollary 2.
Those embeddings already preserve KL polynomials.\\
Principal-lower-interval transfers&Known: Brenti--Caselli--Marietti
\cite[Theorem 7.8 and Corollary 8.4]{bcm} and Delanoy \cite{delanoy}.\\
Small arbitrary intervals&Known through rank six for arbitrary
Coxeter systems: Barkley--Gaetz--Lam \cite[Corollary 1.4]{bgl}.\\
Small finite Weyl/type $A$ intervals&Known through rank ten for Weyl
groups and rank twelve when both realizations have type $A$:
Esposito--Marietti--Stella
\cite[Corollary 1.3(ii)]{ems}.\\
Arbitrary realizations of $I_n$&Theorem~\ref{thm:bad} is a formal
consequence of full invariance and Gern's calculations.
Its unrestricted scope exceeds these short-interval results,
but an independent novelty claim is not established.\\
\bottomrule
\end{tabular}
\end{center}
The rank-eleven $I_6$ test and rank-twenty-one $I_8$ test should not
be described as outside \emph{every} prior result merely because their
ranks exceed ten. A particular realization may be handled by special
matchings, geometry, parabolic cancellation, or pattern embedding.
Nor does the rank-twelve type $A$ statement automatically compare a
type $A$ realization with a type $D$ realization; its precise ambient
hypotheses matter. The defensible claim here is the new unrestricted
\emph{deduction}, conditional on the supplied full theorem.
For $I_6$, the shape certificate above specifically rules out
principal-lower and lattice realizations. It does not rule out
every other older transport mechanism for a particular realization.

\section{Verification, computations, and remaining research}
The project saves the source papers and extracted text, a literature
audit, and executable Python code under \texttt{computations/}.
The finite verification exhaustively compares two implementations on
54,907 comparable pairs (including equal endpoints) and 201,600
ordered Bruhat comparisons in
$A_3,A_4,B_3,D_4,B_4$. They use BFS versus inversion-formula lengths
and subword versus lifting-property Bruhat comparison; both implement
the same standard KL recurrence. The independent $R$ recurrence adds
a different polynomial identity check for the certified examples.
The saved JSON report includes six complete order-isomorphism
certificates and the targeted $D_6$ result. The $D_8$ base computation
is inherited from the thesis and was not independently reproduced.
The executed dependency-free command is
\begin{quote}\ttfamily\small
python3 computations/run\_search.py
\end{quote}
It can also be run through \texttt{uv run --offline} using an installed
Python; project instructions give the exact command and certificate
verification procedure. The exceptional root checks run with
\texttt{python3 research/exceptional\_cosets.py}.
Arithmetic degree tables, small-group KL calculations, and theorem
proofs have different evidentiary roles: the table checks formulas;
small-group calculations check concrete examples; neither proves
unrestricted combinatorial invariance or whole-type coverage.

The logical audit of Theorem~\ref{thm:bad} is short and exact:
\begin{enumerate}
\item The interval certificate supplies the whole closed poset, hence its rank.
\item The external invariance theorem identifies the two ordinary polynomials.
\item Gern's Lemma 2.3.9 permits the $e$-bottom $a$-function degree bound.
\item Proposition 4.5.1 evaluates that source $a$-function.
\item Taking integer floors and the ordinary degree bound yields $B_n$.
\item Lemma 4.5.5 supplies the original computational base values;
our certificate independently reproduces and expands the $n=6$
calculation. The $n=8$ value remains inherited from the thesis and
removes the allowed $q^{10}$ term.
\end{enumerate}
No target $a$-function, target fully commutative status, or
target root-label compatibility is assumed. No assertion is made
for unequal-parameter polynomials, parabolic variants, or
$p$-Kazhdan--Lusztig polynomials.

\paragraph{Update from the expanded search.}
The companion paper \emph{Leading Kazhdan--Lusztig coefficients with
fully commutative lower endpoints in type $E$} gives a
computer-assisted extension to entire finite types $E_6,E_7,E_8$.
Independent parabolic searches certify $E_8$ without enumerating its full
group. The note also proves explicit unbounded full-support upper families
in affine $E_8$ and indefinite $E_{10},E_{13}$. The affine family has zero
$\mu$ for every FC lower endpoint; the indefinite families have the FC-bottom
$0/1$ bound. Classical cell theory and parity supply the structural proof.
These direct results enlarge the conditional abstract model class; they do
not provide new interval isomorphisms. Their historical priority remains
unconfirmed. The expanded search also independently
verifies a fully commutative-bottom affine-$D_4$ pair with $\mu=2$:
with central node $2$ and leaves $0,1,3,4$, put $x=s_0s_1s_3s_4$ and
$w=xs_2x$. Then $P_{x,w}=1+3q+2q^2$. The cube-interval interpretation is
classical, and Mongelli's 2011 preprint \emph{Kazhdan--Lusztig polynomials
of boolean elements}, Corollary 2, explicitly covers the star pattern and
its Catalan leading coefficient. This is a rediscovery, not a novelty claim.

Large $\mu$-values in affine types $B,C,D$ are known
\cite{chen}; they are not counterexamples to an FC-bottom assertion
without checking their lower endpoints. Likewise the initial literature
search's failure to find a complete exceptional-type or affine-$D$
FC-bottom theorem is not a certification that such a theorem remains
open. The separate literature file records queries, exact locators,
version differences, and access limitations, including the university
repository's blocked Chmutov PDF download.

For a substantive next theorem, the missing step is a structural
recognition result supplying FC-bottom $D$ witnesses for a naturally
defined family in exceptional or affine type \emph{beyond} classical
parabolic or subsystem embeddings. Another promising target is a
certified non-parabolic realization of $I_6$ or $I_8$.
Such a high-rank nonclassical realization has not been proved. The exceptional
companion supplies separate direct finite-type and explicit infinite-family
extensions. Its rank-uniform construction also gives full-support upper
families in every $E_{4r+1}$, $r\geq3$, with a second unbounded parameter.
Their FC-bottom leading coefficients are $0$ or $1$, and all non-cover
coefficients vanish. A symbolic $D_{4r}$ sign-change proof establishes
real-root membership; no KL polynomial computation is needed. Historical
priority remains unconfirmed, and full polynomials are not determined.
This note records completed conditional deductions
and concrete classical applications without replacing those gaps by
an unproved novelty claim.

\begin{thebibliography}{99}
\bibitem{gern} T. C. Gern, \emph{Leading Coefficients of
Kazhdan--Lusztig Polynomials in Type $D$}, Ph.D. thesis, University
of Colorado Boulder, 2013. \url{https://arxiv.org/abs/1304.6074}.
\bibitem{ci} OpenAI, \emph{Combinatorial invariance of
Kazhdan--Lusztig polynomials}, preprint, September 24, 2026.
\url{https://github.com/openai/math/blob/main/preprints/Combinatorial-Invariance-of-Kazhdan-Lusztig-Polynomials-September-24-2026/paper.pdf}.
\bibitem{green} R. M. Green, \emph{Leading coefficients of
Kazhdan--Lusztig polynomials and fully commutative elements},
J. Algebraic Combin. \textbf{30} (2009), 165--171.
\url{https://arxiv.org/abs/0801.1650}.
\bibitem{chmutov} M. S. Chmutov, \emph{The structure of $W$-graphs
arising in Kazhdan--Lusztig theory}, Ph.D. thesis, University of
Michigan, 2014, Section 2.3, Theorem 2.3.8.
\url{https://hdl.handle.net/2027.42/108802}.
\bibitem{stembridge} J. R. Stembridge, \emph{On the fully commutative
elements of Coxeter groups}, J. Algebraic Combin. \textbf{5} (1996),
353--385. \url{https://doi.org/10.1007/BF00193185}.
\bibitem{woo} A. Woo, \emph{Interval pattern avoidance for arbitrary
root systems}, Canad. Math. Bull. \textbf{53} (2010), 757--762.
\url{https://arxiv.org/abs/math/0611328}.
\bibitem{bb} A. Bj\"orner and F. Brenti,
\emph{Combinatorics of Coxeter Groups}, Graduate Texts in Mathematics
231, Springer, 2005. \url{https://doi.org/10.1007/3-540-27596-7}.
\bibitem{humphreys} J. E. Humphreys, \emph{Reflection Groups and
Coxeter Groups}, Cambridge Studies in Advanced Mathematics 29,
Cambridge University Press, 1990.
\url{https://doi.org/10.1017/CBO9780511623646}.
\bibitem{lusztig} G. Lusztig, \emph{Cells in affine Weyl groups, II},
J. Algebra \textbf{109} (1987), 536--548.
\url{https://doi.org/10.1016/0021-8693(87)90154-2}.
\bibitem{bcm} F. Brenti, F. Caselli, and M. Marietti,
\emph{Special matchings and Kazhdan--Lusztig polynomials},
Adv. Math. \textbf{202} (2006), 555--601.
\url{https://doi.org/10.1016/j.aim.2005.01.011}.
\bibitem{delanoy} E. Delanoy, \emph{Combinatorial invariance of
Kazhdan--Lusztig polynomials on intervals starting from the identity},
J. Algebraic Combin. \textbf{24} (2006), 437--463.
\url{https://doi.org/10.1007/s10801-006-0014-7}.
\bibitem{bgl} G. T. Barkley, C. Gaetz, and T. Lam,
\emph{Combinatorial invariance for the coefficient of $q$ in
Kazhdan--Lusztig polynomials}, 2026, arXiv:2601.07793v2.
\url{https://arxiv.org/abs/2601.07793v2}.
\bibitem{ems} F. Esposito, M. Marietti, and S. Stella,
\emph{Flip Combinatorial Invariance and Weyl groups}, 2026,
arXiv:2509.16433v2. \url{https://arxiv.org/abs/2509.16433v2}.
\bibitem{dyer} M. Dyer, \emph{On the ``Bruhat graph'' of a Coxeter
system}, Compositio Math. \textbf{78} (1991), 185--191.
\url{https://www.numdam.org/item/CM_1991__78_2_185_0/}.
\bibitem{jonestraces} R. M. Green, \emph{Generalized Jones traces and
Kazhdan--Lusztig bases}, J. Pure Appl. Algebra \textbf{211} (2007),
744--772. \url{https://arxiv.org/abs/math/0509362}.
\bibitem{jones} B. C. Jones, \emph{Leading coefficients of
Kazhdan--Lusztig polynomials for Deodhar elements}, J. Algebraic
Combin. \textbf{29} (2009), 229--260.
\url{https://arxiv.org/abs/0711.1391}.
\bibitem{chen} P. Chen, \emph{Non-trivial Kazhdan--Lusztig
coefficients of affine Weyl groups}, preprint, 2018.
\url{https://arxiv.org/abs/1804.02353}.
\end{thebibliography}
\end{document}
